\documentclass[12pt]{article}
\usepackage[T1]{fontenc}
\usepackage[utf8]{inputenc}
\usepackage{lmodern}
\usepackage{textcomp}
\usepackage{enumitem}
\usepackage{geometry}
\geometry{margin=1in}
\begin{document}
\begin{center}
Answer Key - Constitutional Law - Undergraduate Level Assessment 2 \
\small (Answers are ordered exactly as in the question paper.)
\end{center}

\section*{Multiple Choice}
\begin{enumerate}[label=\arabic*.]
\item \textbf{Answer:} C
\\ \textit{Options:} A) There is no relationship between law and morality.; B) Law is always in advance of moral ideas.; C) The law is inextricably bound up with morals.; D) Morality is generally in advance of the law.
\\ \textit{Note:} Non-positivist legal theorists argue that law cannot be separated from moral principles, asserting that legal systems are fundamentally influenced by societal values and ethical considerations, which makes the law inherently connected to morality.
\item \textbf{Answer:} B
\\ \textit{Options:} A) Because its justification depends on the concept employed.; B) Because any definition of punishment should be value-neutral.; C) Because the concept of punishment is controversial.; D) Because the practice of punishment is separate from its justification.
\\ \textit{Note:} Separating the concept of punishment from its justification is crucial because a value-neutral definition allows for an objective understanding of punishment that is not influenced by subjective moral or ethical beliefs, ensuring fair application under constitutional law.
\item \textbf{Answer:} C
\\ \textit{Options:} A) In hard cases judges generally decide cases on the basis of rights.; B) The rights of the parties feature in the determination of most cases before the courts.; C) Judges exercise strong discretion.; D) Judges seek the best 'fit' with constitutional and institutional history.
\\ \textit{Note:} The proposition that judges exercise strong discretion is inconsistent with Dworkin's claim because it implies that multiple interpretations and outcomes are possible in legal questions, contradicting the idea that there is only one right answer.
\item \textbf{Answer:} A
\\ \textit{Options:} A) The POP will choose wealth over a compassionate society.; B) The POP will choose equality over power.; C) The POP will be unselfish.; D) The POP will choose to protect the disabled.
\\ \textit{Note:} The correct answer reflects Rawls' assertion that individuals behind the veil of ignorance would prioritize social primary goods, such as wealth and resources, over subjective concepts like compassion, as these goods are essential for ensuring fairness and enabling individuals to pursue their own visions of the good life.
\item \textbf{Answer:} D
\\ \textit{Options:} A) Incorporationism' accepts that judges may decide cases by reference to moral factors.'; B) A soft positivist rejects the role of morality in the description of law.; C) Sometimes the definition of law includes moral considerations.; D) Judges lack strong discretion.
\\ \textit{Note:} The correct answer is based on the understanding that if soft positivism allows for the inclusion of moral criteria in the rule of recognition, it implies that judges have some degree of discretion to interpret the law in light of those moral considerations, thus making the proposition that judges lack strong discretion inconsistent with this view.
\end{enumerate}

\section*{True / False}
\begin{enumerate}[label=\arabic*.]
\setcounter{enumi}{5}
\item \textbf{Answer:} False
\\ \textit{Options:} A) True; B) False
\\ \textit{Note:} Under the formalist approach to legal reasoning, legislative acts are considered the primary source of law, with judicial decisions primarily serving to interpret and apply those legislative provisions rather than supersede them.
\item \textbf{Answer:} True
\\ \textit{Options:} A) True; B) False
\\ \textit{Note:} The Sociological School of jurisprudence posits that law is a tool for social change and must reflect the needs and values of society to effectively promote justice and achieve social goals.
\item \textbf{Answer:} True
\\ \textit{Options:} A) True; B) False
\\ \textit{Note:} Kelsen's pure theory of law focuses solely on the structure and function of legal norms without considering external factors, making it a jurisprudential analysis distinct from other legal theories.
\item \textbf{Answer:} True
\\ \textit{Options:} A) True; B) False
\\ \textit{Note:} The Analytical School emphasizes the role of judicial recognition in establishing customs as law, asserting that a custom gains legal validity only when it is acknowledged and enforced by the courts.
\item \textbf{Answer:} False
\\ \textit{Options:} A) True; B) False
\\ \textit{Note:} Nozick argues against redistribution of wealth by asserting that individuals have the right to their holdings as long as they were acquired justly, emphasizing that any enforced redistribution would violate individual rights and thus be unjust.
\end{enumerate}

\section*{Long Form Response}
\begin{enumerate}[label=\arabic*.]
\setcounter{enumi}{10}
\item \textbf{Answer:} The positive obligations of the state as established by the European Court of Human Rights (ECHR) primarily focus on ensuring the protection of individuals' rights, particularly under the European Convention on Human Rights. While states are required to take proactive measures to safeguard rights such as the right to life and the prohibition of inhumane treatment, the obligation to provide housing for all homeless individuals is not explicitly mandated. This exclusion stems from the understanding that while states must address homelessness and implement effective policies, they are not necessarily required to guarantee housing as a fundamental right. Instead, the ECHR emphasizes the need for states to create conditions that prevent homelessness and to support individuals in achieving housing stability, thereby balancing state resources and individual rights.
\\ \textit{Note:} correct because it accurately identifies that the ECHR's positive obligations require states to protect human rights but do not extend to an explicit duty to provide housing for all homeless individuals, reflecting a distinction between safeguarding rights and ensuring the provision of social services.
\item \textbf{Answer:} Conciliation and mediation are both alternative dispute resolution processes, but they differ significantly in their execution and roles of participants. Mediation typically involves a mediator appointed by the consent of the disputing parties, who facilitates dialogue and encourages mutual agreement without imposing any solutions. In contrast, conciliation is often carried out by a commission that conducts an impartial examination of the dispute and proposes specific settlement terms for the parties to consider. The implications for constitutional law arise from the differing levels of authority and involvement; while mediation emphasizes voluntary solutions, conciliation's structured approach can lead to more definitive resolutions that may influence legal interpretations and enforceability in constitutional contexts. These distinctions highlight the flexibility of dispute resolution mechanisms and their potential impact on upholding constitutional principles.
\\ \textit{Note:} The answer correctly highlights that mediation focuses on facilitating dialogue and mutual agreement without imposing solutions, while conciliation involves an impartial examination of the dispute with proposed settlement terms, illustrating their distinct processes and participant roles in dispute resolution.
\end{enumerate}

\vfill
\noindent\textit{End of Answer Key}
\end{document}